\documentclass[12pt,letterpaper]{article}

% --- Paquetes Principales ---
\usepackage[utf8]{inputenc}
\usepackage[spanish,es-tabla]{babel}
\usepackage[margin=2.5cm]{geometry}
\usepackage{graphicx}
\usepackage{amsmath,amssymb}
\usepackage{booktabs}
\usepackage{float}
\usepackage{tabularx}
\usepackage{array}
\usepackage{longtable}
\usepackage{fancyhdr}
\usepackage{titlesec}
\usepackage{microtype}
\usepackage{xcolor}
\usepackage[hidelinks,pdfencoding=auto,psdextra]{hyperref}

% --- Colores ---
\definecolor{primary}{RGB}{20, 50, 90}
\definecolor{secondary}{RGB}{70, 110, 150}

% --- Configuración de Encabezados y Pies de Página ---
\pagestyle{fancy}
\fancyhf{}
\fancyhead[L]{\small\scshape Sistemas de Información}
\fancyhead[R]{\small\scshape [Título del Documento]}
\fancyfoot[C]{\thepage}
\renewcommand{\headrulewidth}{0.4pt}
\renewcommand{\footrulewidth}{0.2pt}

% --- Formato de Secciones y Subtítulos ---
\titleformat{\section}
  {\normalfont\Large\bfseries\color{primary}}{\thesection}{1em}{}
\titleformat{\subsection}
  {\normalfont\large\bfseries\color{secondary}}{\thesubsection}{1em}{}
\titleformat{\subsubsection}
  {\normalfont\normalsize\bfseries}{\thesubsubsection}{1em}{}

\setlength{\parskip}{0.6em}
\setlength{\parindent}{0pt}

\begin{document}

% ==============================================================================
% PORTADA
% ==============================================================================

\begin{titlepage}
    \centering
    
    % Logo Institucional
    \includegraphics[width=0.3\textwidth]{unal_logo.png}\par
    \vspace{0.8cm}
    % Materia e Institución
    {\scshape\Large Universidad Nacional de Colombia\par}
    {\scshape\large Facultad de Ingeniería\par}
    {\scshape\normalsize Departamento de Ingeniería de Sistemas e Industrial\par}
    \vspace{0.5cm}
    {\Large\textbf{Sistemas de Información}\par}
    
 
    
    % Título y Subtítulo
    {\huge\bfseries\color{primary} Hito 1\par}
    \vspace{0.4cm}
    

    
    % Integrantes (7 miembros)
    {\large
    Sebastian Fernando Gutiérrez Rojas\par
    Ronald Arturo Chávez Cabrera\par
    Nombre Completo Integrante 3\par
    Nombre Completo Integrante 4\par
    Nombre Completo Integrante 5\par
    Nombre Completo Integrante 6\par
    Nombre Completo Integrante 7\par
    }
    
    \vfill
    
    % Docente y Fecha
    {\normalsize\textbf{Docente:} Jhon Alexander Garcia Camargo}\par
    \vspace{0.4cm}
    {\normalsize Bogotá D.C., Colombia\par}
    {\normalsize 2026}
\end{titlepage}

% ==============================================================================
% TABLA DE CONTENIDO
% ==============================================================================
\newpage
\tableofcontents
\newpage
\section{Diagnóstico y alineación estratégica}
\subsection{Análisis de las 5 Fuerzas de Porter}
\subsection{Mapeo de la Cadena de Valor identificando puntos de dolor tecnológicos.}
\subsection{Definición inicial del problema bajo el estándar de Diseño de Ingeniería.}
\subsubsection{Identificación de necesidades}

La definición de necesidades se aborda siguiendo las fases clásicas de la ingeniería de requerimientos (ISO/IEC/IEEE 29148): \textbf{elicitación}, \textbf{análisis y negociación}, \textbf{especificación} y \textbf{validación}. Este apartado cubre las dos primeras y deja preparada la trazabilidad hacia la especificación formal de requerimientos funcionales y no funcionales.

\paragraph{Enfoque y metodología.}

\begin{table}[H]
\centering
\small
\caption{Fases de ingeniería de requerimientos aplicadas.}
\begin{tabularx}{\textwidth}{@{}l X X@{}}
\toprule
\textbf{Fase} & \textbf{Actividad realizada} & \textbf{Producto} \\
\midrule
Elicitación & Análisis documental del repositorio actual (Dropbox, más de 23.000 archivos), análisis de las 5 Fuerzas de Porter y mapeo de la cadena de valor con sus puntos de dolor. & Lista bruta de problemas y evidencias. \\
Análisis & Agrupación por dominio, identificación de interesados, formulación de necesidades y priorización. & Catálogo de necesidades \texttt{NEC-xx}. \\
Negociación & Delimitación del alcance (área financiera excluida) y definición de restricciones y supuestos. & Alcance acordado y lista de preguntas por validar. \\
Especificación & Derivar requerimientos funcionales (\texttt{RF}) y no funcionales (\texttt{RNF}) verificables desde cada necesidad. & Trazabilidad NEC $\rightarrow$ RF/RNF. \\
\bottomrule
\end{tabularx}
\end{table}

 Una \emph{necesidad} describe qué requiere la organización para resolver un problema, sin prescribir la solución. Un \emph{requerimiento} es la traducción verificable de esa necesidad en una capacidad del sistema. 

\paragraph{Formalización del problema de ingeniería.}

Aplicando el estándar de Diseño en Ingeniería, el problema se formaliza en la siguiente tabla:

\begin{table}[H]
\centering
\small
\caption{Formalización del problema bajo el estándar de Diseño en Ingeniería.}
\begin{tabularx}{\textwidth}{@{}l X@{}}
\toprule
\textbf{Elemento} & \textbf{Enunciado} \\
\midrule
\textbf{Situación actual} & Loard Construcciones opera con Dropbox como gestor documental \emph{de facto}: más de 23.000 archivos en estructura jerárquica, sin metadatos, con versionado manual (\texttt{v1}, \texttt{v6}), copias en conflicto y sin \emph{audit trail}. La operación de campo reporta por WhatsApp y la consolidación físico-financiera se realiza al cierre del proyecto. \\
\textbf{Problema} & La arquitectura de información actual \textbf{no permite} trazabilidad en tiempo real, ni control de versiones confiable, ni explotación analítica de datos históricos (APU, rendimientos, consumibles), lo que se traduce en: (i) erosión del margen por desviaciones detectadas tardíamente, (ii) riesgo legal y contractual por vencimientos no controlados, (iii) reprocesos administrativos en la consolidación del dossier de calidad. \\
\textbf{Necesidad} & La organización requiere un \textbf{ERP complementario} que digitalice la gestión operativa de campo y taller, la gestión comercial (APU y cotizaciones) y el control documental del SIG, integrando captura móvil, control de versiones, metadatos y flujos de aprobación con firma digital. \\
\textbf{Restricción} & El área financiera y contable queda explícitamente \textbf{fuera del alcance}. Se debe cumplir con ISO 9001:2015, SG-SST (Decreto 1072 de 2015), RUC y protección de datos (Ley 1581 de 2012). \\
\textbf{Criterio de éxito} & Reducción \(\geq 50\%\) del tiempo administrativo de consolidación documental; \(\geq 95\%\) de los reportes diarios capturados digitalmente; cero no conformidades por uso de versiones obsoletas en auditoría externa; 100\% de los vencimientos SST con alerta anticipada \(\geq 30\) días. \\
\bottomrule
\end{tabularx}
\end{table}

\begin{table}[H]
\centering
\small
\caption{Relación causa--efecto del diagnóstico.}
\begin{tabularx}{\textwidth}{@{}l X X@{}}
\toprule
\textbf{Síntoma observable} & \textbf{Causa raíz} & \textbf{Efecto en el negocio} \\
\midrule
``Falta dinero'' al cierre del proyecto & Sin trazabilidad de consumibles en tiempo real (NEC-02) & Margen erosionado sin capacidad de corrección \\
Informes diarios por WhatsApp & Ausencia de canal digital estructurado en campo (NEC-01) & Avance físico-financiero desfasado \\
Versiones obsoletas en auditoría & Versionado manual en el nombre del archivo (NEC-12) & No conformidades ISO 9001; riesgo contractual \\
Copias en conflicto de \texttt{GT-FR25} & Colisión de sincronización en Dropbox (NEC-13) & Pérdida de evidencia de pruebas hidrostáticas \\
Cotizaciones tardías & APU heterogéneos y búsqueda manual de históricos (NEC-06) & Pérdida de licitaciones por sobretiempo \\
Vencimientos SST no detectados & Control manual en hojas de cálculo (NEC-18) & Riesgo legal (Decreto 1072 de 2015) y de seguridad \\
Dossier de calidad tardío & Transcripción manual de ensayos NDT (NEC-05) & Retraso en radicación de actas y facturación \\
\bottomrule
\end{tabularx}
\end{table}

\paragraph{Partes interesadas (stakeholders).}

La matriz clasifica a los interesados según su \textbf{poder} e \textbf{interés} , y deriva la estrategia de gestión del cuadrante resultante.

\begin{table}[H]
\centering
\small
\caption{Matriz poder--interés de interesados.}
\begin{tabularx}{\textwidth}{@{}l c c X@{}}
\toprule
\textbf{Interesado} & \textbf{Poder} & \textbf{Interés} & \textbf{Estrategia de gestión} \\
\midrule
Gerencia general & Alto & Alto & \textbf{Gestionar de cerca.} Patrocinador del proyecto; validación de alcance y presupuesto. \\
Clientes (Chevron, Primax) & Alto & Bajo & \textbf{Mantener satisfechos.} Auditorías y requisitos contractuales condicionan el diseño. \\
Coordinación del SIG / calidad & Alto & Alto & \textbf{Gestionar de cerca.} Dueño funcional de NEC-12 a NEC-17. \\
Gestión Humana y SST & Alto & Alto & \textbf{Gestionar de cerca.} Dueño funcional de NEC-18 y NEC-19; riesgo legal directo. \\
Ingenieros residentes y supervisores & Medio & Alto & \textbf{Mantener informados.} Usuarios clave de captura en campo. \\
Área comercial y de costos & Medio & Alto & \textbf{Mantener informados.} Usuarios clave de NEC-06 y NEC-07. \\
Almacén y logística & Bajo & Alto & \textbf{Mantener informados.} Usuarios de NEC-02 y NEC-03. \\
Compras & Bajo & Medio & \textbf{Monitorear.} Usuarios de NEC-08. \\
Cuadrillas y personal de campo & Bajo & Alto & \textbf{Mantener informados.} Usabilidad móvil crítica. \\
Área financiera y contable & Alto & Nulo & \textbf{Monitorear.} Explícitamente fuera del alcance; riesgo de fricción política. \\
\bottomrule
\end{tabularx}
\end{table}

\paragraph{Catálogo de necesidades identificadas.}

Cada necesidad se formula bajo la estructura ``La empresa necesita [capacidad] para [beneficio]'' y se vincula con su evidencia en el diagnóstico. La prioridad usa MoSCoW: \textbf{M} (debe), \textbf{S} (debería), \textbf{C} (podría).

\medskip
\noindent\textbf{A. Operación de campo y taller}

\begin{longtable}{@{}p{1.4cm} p{5.6cm} p{4.6cm} p{1.4cm}@{}}
\toprule
\textbf{ID} & \textbf{Necesidad} & \textbf{Evidencia} & \textbf{Prior.} \\
\midrule
\endhead
NEC-01 & Registrar de forma estructurada el informe diario de actividades de cada cuadrilla, para sustituir los reportes por WhatsApp y consolidar el avance sin retrasos. & Informalidad en reportes de campo; formato \texttt{GT-FR48} es soporte de facturación (Cadena de valor, Operaciones). & M \\
NEC-02 & Controlar en tiempo real el consumo de consumibles (electrodos, discos, gases) por frente de obra, para detectar desviaciones antes del cierre del proyecto. & Carencia de trazabilidad; la conciliación se hace al cierre y se descubre que ``falta dinero'' (Logística interna). & M \\
NEC-03 & Contar con hoja de vida, estado, fecha de adquisición y ubicación de cada herramienta y equipo, para reducir pérdidas y asignarlos por frente. & Inventario estático y ciego; descontrol en devoluciones; sin stock de seguridad (Logística interna). & S \\
NEC-04 & Medir el avance físico y la productividad real de las cuadrillas frente a lo cotizado, para evaluar la rentabilidad de contratos marco y proyectos. & Rendimientos sin analítica; desfase físico-financiero (Operaciones). & M \\
NEC-05 & Consolidar los resultados de ensayos (NDT, hidrostáticas, cilindros de concreto) y el libro de tubería en un dossier de calidad, sin transcripción manual a hojas de cálculo. & Consolidación manual del dossier; \texttt{GT-FR64} (Weld Book) y formatos de ensayos (Logística externa; Gestión Técnica). & M \\
\bottomrule
\end{longtable}

\noindent\textbf{B. Gestión comercial y gerencial}

\begin{longtable}{@{}p{1.4cm} p{5.6cm} p{4.6cm} p{1.4cm}@{}}
\toprule
\textbf{ID} & \textbf{Necesidad} & \textbf{Evidencia} & \textbf{Prior.} \\
\midrule
\endhead
NEC-06 & Estandarizar el Análisis de Precios Unitarios (APU) y conservar un histórico consultable de tarifas y rendimientos, para cotizar con consistencia y rapidez. & Plantillas Excel heterogéneas; búsquedas en carpetas de hace años o por llamadas (Marketing y ventas). & M \\
NEC-07 & Estandarizar la generación de cotizaciones (con y sin AIU) y su seguimiento, para evitar sobretiempos y reprocesos al licitar. & Sobreesfuerzo operativo en cotización; formato \texttt{GC-FR04} (Gestión Comercial). & S \\
NEC-08 & Visualizar el ciclo completo de una compra (requisición, cotización, aprobación, despacho y entrega), para reducir compras urgentes con sobreprecio. & Falta de trazabilidad del lead time; autorizaciones sin flujo formal; sobrecostos en Buenaventura (Compras). & M \\
NEC-09 & Disponer de un tablero gerencial operativo con avance físico, costos directos de obra y alertas por proyecto, sin depender de revisiones manuales ni llamadas. & Opacidad de costos directos; facturación ciega; desfase de información (Infraestructura). & M \\
NEC-10 & Analizar la concentración de ingresos por cliente y contrato, y la desviación entre lo cotizado y lo ejecutado, para sustentar el margen de imprevistos (hoy 7--10\% empírico). & 80\% de la facturación en un solo cliente (Porter, poder de los clientes); contingencias no analizadas (Postventa). & S \\
NEC-11 & Mantener un repositorio de lecciones aprendidas y fallas por proyecto que alimente nuevas cotizaciones. & Ausencia de lecciones aprendidas (Postventa). & C \\
\bottomrule
\end{longtable}

\noindent\textbf{C. Gestión documental del Sistema Integrado de Gestión (SIG)}

\begin{longtable}{@{}p{1.4cm} p{5.6cm} p{4.6cm} p{1.4cm}@{}}
\toprule
\textbf{ID} & \textbf{Necesidad} & \textbf{Evidencia} & \textbf{Prior.} \\
\midrule
\endhead
NEC-12 & Controlar versiones de forma automática y garantizar que el personal use solo la versión vigente, con las obsoletas archivadas. & Versionado manual en el nombre (\texttt{v1}, \texttt{v6}); carpeta \texttt{Documentos Obsoletos}; copias locales divergentes (Diagnóstico Dropbox; Desarrollo tecnológico). & M \\
NEC-13 & Evitar la edición simultánea conflictiva sobre un mismo documento, para no perder información de pruebas ni informes. & Archivos ``copia en conflicto'' (p.\,ej. \texttt{GT-FR25}). & M \\
NEC-14 & Buscar documentos por atributos (cliente, sede, año, proceso, tipo), no solo por ruta de carpetas. & Estructura exclusivamente en carpetas anidadas, sin metadatos. & M \\
NEC-15 & Restringir el acceso por rol a información sensible (nómina, salud ocupacional, tarifas y márgenes). & Permisos a nivel de carpeta; riesgo de alteración o eliminación. & M \\
NEC-16 & Aprobar y firmar digitalmente actas y autorizaciones, con validez probatoria y sello de tiempo, evitando imprimir y escanear. & Firmas manuscritas escaneadas en PDF. & S \\
NEC-17 & Registrar quién consultó, modificó o descargó cada documento, para sustentar auditorías de clientes y de certificación. & Sin audit trail en el gestor actual. & M \\
\bottomrule
\end{longtable}

\noindent\textbf{D. Talento humano y seguridad y salud en el trabajo}

\begin{longtable}{@{}p{1.4cm} p{5.6cm} p{4.6cm} p{1.4cm}@{}}
\toprule
\textbf{ID} & \textbf{Necesidad} & \textbf{Evidencia} & \textbf{Prior.} \\
\midrule
\endhead
NEC-18 & Recibir alertas anticipadas de vencimiento de cursos (alturas, espacios confinados), calificaciones de soldadores y exámenes médicos ocupacionales. & Riesgo operativo y legal por vencimientos inadvertidos (Gestión de RR.\,HH.; matriz de competencias). & M \\
NEC-19 & Digitalizar el registro de inducciones y capacitaciones hechas en campo, para no depender de soportes físicos que viajan entre ciudades. & Dispersión documental y riesgo de pérdida ante auditorías del cliente. & S \\
\bottomrule
\end{longtable}

\paragraph{Análisis y priorización.}

La priorización combina tres criterios: (i) riesgo legal, contractual o de seguridad si la necesidad no se atiende, (ii) impacto económico (margen, sobrecostos, tiempo administrativo) y (iii) frecuencia con la que ocurre el problema.

\begin{table}[H]
\centering
\small
\caption{Resumen de priorización MoSCoW.}
\begin{tabularx}{\textwidth}{@{}l X X@{}}
\toprule
\textbf{Prioridad} & \textbf{Necesidades} & \textbf{Justificación} \\
\midrule
Debe (M) & NEC-01, 02, 04, 05, 06, 08, 09, 12, 13, 14, 15, 17, 18 & Afectan directamente el margen, la seguridad de las personas, la trazabilidad exigida por clientes petroleros o la conformidad ante auditorías (ISO 9001, RUC). \\
Debería (S) & NEC-03, 07, 10, 16, 19, 20 & Mejoran eficiencia y control, pero el negocio puede operar sin ellas en una primera versión. \\
Podría (C) & NEC-11 & Valor estratégico a largo plazo; depende de contar antes con datos históricos limpios. \\
\bottomrule
\end{tabularx}
\end{table}

\textbf{Observaciones del análisis.}
\begin{itemize}
    \item Las necesidades NEC-01, 02, 04 y 09 forman una cadena: sin captura diaria estructurada (NEC-01) no hay consumo (NEC-02), productividad (NEC-04) ni tablero gerencial (NEC-09). Conviene abordarlas en el mismo incremento.
    \item Las necesidades del dominio C (NEC-12 a 17) son las que justifican el gestor documental, y resuelven los seis puntos de dolor del diagnóstico de Dropbox (concurrencia, versiones, metadatos, RBAC, firmas y auditoría).
    \item NEC-18 es la de mayor riesgo legal inmediato ya que trabajo en alturas y espacios confinados, Decreto 1072 de 2015, y tiene bajo costo de implementación; es candidata a primera entrega.
\end{itemize}

\paragraph{Restricciones y supuestos.}

\begin{itemize}
    \item \textbf{Alcance:} el área financiera y contable no participa; los sistemas contables, balances y conciliaciones tributarias quedan fuera. El ERP cubre la gestión operativa de campo y taller y la gestión gerencial y comercial.
    \item \textbf{Normativas:} el sistema debe respaldar la información documentada de ISO 9001:2015, los requisitos del SG-SST (Decreto 1072 de 2015) y del RUC. Para datos de personal aplica la protección de datos personales (Ley 1581 de 2012) y, para firmas digitales, la Ley 527 de 1999.
    \item \textbf{Contexto de uso:} usuarios en campo con distintos niveles de familiaridad tecnológica, en cinco sedes (Cartagena, Bogotá, Medellín, Yumbo y Buenaventura).
    \item \textbf{Supuesto:} el repositorio actual se conserva como fuente de migración (más de 23.000 archivos y la codificación \texttt{[Proceso]-[Tipo][Número] v[Versión]}).
\end{itemize}

\subsection{Modelo de Alineación Estratégica (SAM)}

\subsection{Análisis de Gobernanza y Marco de Decisión}

\subsection{Balanced Scorecard: Definición del Problema mediante KPI y KRI}




\section{Discovery de producto}
\subsection{Opportunity Solution Tree (OST)}
\subsection{Customer Journey Map y Definición de Roles de Usuario}



\section{Ingeniería de procesos (BPMN 2.0)}
\subsection{Diagrama As-Is: Mapa de Calor y Cuellos de Botella Actuales}
\subsection{Diagrama To-Be: Proceso Optimizado con el Sistema de Información}

\section{Especificación de requerimientos}
\subsection{Product Backlog: Historias de Usuario bajo Criterios INVEST y Priorización MoSCoW}
\subsection{Matriz de Requerimientos No Funcionales (RNF) y Restricciones Técnicas y Legales (Ley 1581)}
\end{document}
